\section{From a global extremizer to a conditional comparison}\label{loc:section}

Fix an actual integer $26\le n\le30$ and assume the unrestricted
predecessor bound $C_{n-1}\le\CE$. Suppose $C=C_n\ge\CE$, and take
the positive closed maximizing pair $(P,H)$ of
Lemma~\ref{gi:extremizer}, with $1\le\rho\le6/5$.
All quantities below belong to this one pair. Write
$\beta=\E[X|X|]$, $\gamma=\E X^3$, $M_j=\E|X|^j$,
$z=H/\sqrt n$ and $\varphi=\phi(z)$.
The moment-defect construction and its original finite subdivisions are
due to Rodziewicz--Ko\l{}odziejek~\cite{RK2026}.
We spell out the continuous implications needed to interpret the finite
rows; the arithmetic verification of a row is a separate obligation.

First, the predecessor bound has a pointwise support consequence.
Set $s_n=\sqrt{n/(n-1)}$, $e_n=n(s_n-1)-1/2$, and
\[
 \mathcal H_n(H,x)=n^{3/2}
 \{\Phi((H-x)/\sqrt{n-1})-\Phi(H/\sqrt n)\}
 +nx\varphi+\tfrac H2(x^2-1)\varphi.
\]
Multiplying \eqref{gi:contact} by $\sqrt n$ and using
$F_{n-1}(H-x)\le\Phi((H-x)/\sqrt{n-1})+
C_{n-1}\rho/\sqrt{n-1}$ gives
\begin{equation}\label{loc:support-test}
 \sqrt n\,\mathcal D(x)\le\mathcal H_n(H,x)
 +C\{-|x|^3+\tfrac32\rho x^2+3\beta x+\rho e_n\}.
\end{equation}
The otherwise present allowance
$n\rho s_n(C_{n-1}-C)_+$ is zero. Since $\mathcal D=0$ on the
full support, a strictly negative upper bound for the right side on
$x\le-R_-$ or $x\ge R_+$ excludes that entire ray, including atoms.
The bound must hold for every real $H$ and every moment and score value
in the incoming intervals. Thus each support refinement uses the actual
integer $n$ and previously established moment intervals. In particular,
a support calculation at $25$ cannot simply be relabelled at $n$.
For the initial radius no such calculation is needed:
Proposition~\ref{gc:initial-support} gives $12/5$ when $\rho\le11/10$
and $13/5$ when $\rho\le6/5$, by one Peano-kernel argument for all
$n\ge25$. This precedes all higher-moment and defect refinements.

Here is the physical resource behind the refinement. Put
$\sigma=\sqrt{10}$ and $\tau(t)=\E[X\mathbf1_{\{X\ge t\}}]$.
For every bounded component $g=(a,b)$ of the complement of the closed
support, $\tau$ is constant there; call the value $\tau_g$. Define
\begin{align*}
 J&=\sum_g\tau_g(b-a)^2,&
 J_1&=\sum_g\tau_g\int_a^b(t-a)t\,dt,\\
 J_2&=\sum_g\tau_g\int_a^b(t-a)t|t|\,dt,&
 D&=(3+\sigma)\rho-\gamma-3J.
\end{align*}
All these expressions are finite by the bounded support already proved.
For $\nu=\mathcal L(H-X_1-\cdots-X_{n-1})$, also put
\[
 \mathcal L=\sum_g\tau_g\int_{[a,b)}(t-a)\,\nu(dt)\ge0.
\]
The half-open interval specifies the convention at atoms.

\begin{lemma}[The physical defect]\label{loc:defect}
One has $D\ge0$. If $D<(\sigma-1)\rho$, there is an actual central
gap $(-A_-,A_+)$ with $A_-,A_+>0$. Put
\[
 p=\PP(X>0),\quad q=1-p,\quad v=pq,\quad
 \mu=\E X_+=\E X_-,\quad
 \bar A_-=\mu/q,\quad\bar A_+=\mu/p,
\]
and $\Delta_\pm=\bar A_\pm-A_\pm\ge0$.
Let $\theta_\pm=\E[(X-\E[X\mid\operatorname{sgn}X])^2;
\ \pm X>0]$ and $\theta=\theta_-+\theta_+$.
Then, with $p_*=(4-\sigma)/2$,
\begin{align}\label{loc:joint-defect}
 D\ge{}&\frac{(6+2\sigma)(1-\theta)^{3/2}}{\sqrt v}
                    (p-p_*)^2\notag\\
 &+6\mu\{(A_-+A_+)(\Delta_-+\Delta_+)+\Delta_-\Delta_+\}
 +3(\sigma+3)A_-\theta_-+3(\sigma+1)A_+\theta_+.
\end{align}
\end{lemma}
\begin{proof}
If zero lies in the support, all gaps have one sign and
$J\le2\int|t|\tau(t)\,dt=\rho$. Hence
$D\ge\sigma\rho-\gamma\ge(\sigma-1)\rho$.
Otherwise let $J_o=J-\mu(A_-+A_+)^2$ and
\[
 T_\pm=\E X_\pm^3-\mu A_\pm^2,
 \qquad R_\pm=\E[X_\pm(X_\pm-A_\pm)].
\]
On a positive outer gap $(a,b)$,
$(b-a)^2\le2\int_a^b(t-A_+)dt$.
Summing and applying Tonelli, and then reflecting for the negative
side, proves $J_o\le T_++T_--2A_+R_+-2A_-R_-$.
Expansion gives the nonnegative decomposition
\[
 D=\mu(\sqrt{\sigma+1}A_--\sqrt{\sigma-1}A_+)^2
 +(\sigma-1)T_++(\sigma+1)T_-+3(T_++T_--J_o).
\]
For the stronger assertion, conditional centering gives
\[
 \E X_\pm^3=\mu\bar A_\pm^2+
 \E[(X_\pm+2\bar A_\pm)(X_\pm-\bar A_\pm)^2;
                         \ \pm X>0]
 \ge\mu\bar A_\pm^2+3A_\pm\theta_\pm,
\]
and $R_\pm=\theta_\pm+\mu\Delta_\pm$.
Insert these inequalities into the preceding outer-gap bound.
Use $\mu^2=(1-\theta)pq$ and
$(\sigma+1)p^2+(\sigma-1)q^2-6pq
=(6+2\sigma)(p-p_*)^2$ to obtain \eqref{loc:joint-defect}.
\end{proof}

The upper bound for this same $D$ comes from contact and Fourier
inversion, rather than from an assumed two-point approximation.
For $w=H^2/(2n)$ the exact conserved identity is
\begin{equation}\label{loc:conserved}
 n(C-\CE)\rho+n\CE\rho(1-e^{-w})
 +\frac{n\phi(0)e^{-w}}6D+n^{3/2}\mathcal L=\mathcal R,
\end{equation}
where, setting
\[
 K_1=\E[X|X|^3]-3\beta-\tfrac32\rho\gamma,\quad
 K_q=M_5/4-\rho M_4/2-\tfrac32\beta\gamma+\rho/2,
 \quad U=J_1+(M_4-3)/6,
\]
the right side is
\begin{equation}\label{loc:raw}
 \mathcal R=C(3\beta J/2+3\rho J_1-3J_2-K_q)
 +H(CK_1+\varphi U)-H^2(C\rho+\gamma\varphi/2)+E_n.
\end{equation}
The remainder has the exact representation
\begin{equation}\label{loc:fourier-remainder}
 E_n=\frac{2\sqrt n}{\pi}\int_0^\infty
 \frac{\Im\{e^{-itH}[f(t)^n-e^{-nt^2/2}(1-in\gamma t^3/6)]\}}
 {t^3}\,dt,\qquad f(t)=\E e^{itX}.
\end{equation}
This integral is absolutely convergent.
To see the identity, write $\nu([x,\infty))\le\mathfrak q(x)$ for
the cubic majorant furnished by \eqref{gi:contact}, with equality on
the support. Evaluate
$\int\tau(x)[\mathfrak q(x)-\nu([x,\infty))]dx$ first gap by gap
and then using the density $\tau$ of the zero-bias law.
The first evaluation is
$\varphi J/(2\sqrt n)+G/n^{3/2}-\mathcal L$, where
$G=3C\beta J/2+(H\varphi+3C\rho)J_1-3CJ_2$.
For the second, contact multiplied by $X$ supplies the first truncated
moment; the payoff $(x-z)|x-z|$ supplies the second truncated moment.
Their difference gives
\begin{align*}
 C\rho+\sqrt n\mathcal L
 &=\frac{\varphi}{6}(\gamma+3J)+(G-P_0)/n+E_n/n,\\
 P_0&=H^2(C\rho+\gamma\varphi/2)
      -H\{CK_1+\varphi(M_4-3)/6\}+CK_q.
\end{align*}
Substitute $\gamma+3J=(3+\sigma)\rho-D$ to obtain
\eqref{loc:conserved}--\eqref{loc:raw}.
Formula~\eqref{loc:fourier-remainder} follows from
$\int_0^\infty(ty-\sin ty)t^{-3}dt=\pi y|y|/4$ by compensated
Fubini; bounded support controls the origin and $t^{-3}$ controls
infinity. No discontinuous distribution function is inverted.

There is a necessary order in using this identity. A first bound
$\mathcal R<r<n c_0$, with $C\rho\ge c_0>0$, proves
$\varphi(\gamma+3J)/6=C\rho+\sqrt n\mathcal L-\mathcal R/n>0$.
Only then may the Gaussian weight be removed, giving
\begin{equation}\label{loc:unweighted}
 n(C-\CE)\rho+\kappa nD+n^{3/2}\mathcal L\le\mathcal R,
 \qquad \kappa=997/15000<\phi(0)/6.
\end{equation}
Every later raw bound $\mathcal R\le r'$ also gives
$e^{-w}\ge1-5r'/(2n)$ and $nD\le r'/\kappa$.
The initial pass therefore requires no defect bound of its own output.
The actual central gap is used only after
$D<(\sigma-1)\rho$ has been established.

On a closed moment cell, the finite calculation bounds
\eqref{loc:fourier-remainder} by a coherent expression
$K_F+s_F|H|$, and bounds the coefficients of \eqref{loc:raw}.
It thereby proves an envelope
$\mathcal R\le\max_j(A_j+F_{j,\pm}|H|-Q_j^{\rm raw}H^2)$,
where every $Q_j^{\rm raw}>0$.
Square completion bounds the envelope on the whole half-line.
The signed curvature uses the actual $\gamma$: if
$\gamma\ge g_0$ and $L_\phi\le\varphi\le U_\phi$, its valid price is
$C\rho+\tfrac12\min(L_\phi g_0,U_\phi g_0)$.
Reflecting an auxiliary Fourier integral does not change this sign.
Keeping the phase cost in \eqref{loc:conserved} further gives
\begin{equation}\label{loc:phase-budget}
 \kappa nD\le\max_j(A_j+F_{j,\pm}|H|-Q_jH^2).
\end{equation}
For example, a previously proved $\mathcal R\le r$ and $D\le d$
allow the additional curvature
$(1-5r/(2n))( .4095\rho-d/15)/2$, when positive.
It is included once in $Q_j$.
Thus a bound on the right side of \eqref{loc:phase-budget} is converted
to physical $D$ using $15000/(997n)$ at the actual integer.

We now give the canonical parameters required by
Section~\ref{mc:section}. Once the central gap is licensed, put
\begin{equation}\label{loc:canonical}
 B=\mathbf1_{\{X>0\}},\qquad
 \ell=\bar A_-+\bar A_+=\mu/v,\qquad
 U=X/\ell+p-B.
\end{equation}
Then $\E[U\mid B]=0$ and $X=\ell(B+U-p)$ exactly.
With $s_i=\E[U^2\mid B=i]$, $s=qs_0+ps_1$ and $\lambda_n=ns$,
\begin{equation}\label{loc:canonical-identities}
 \ell^2=\frac{1-\theta}{v}=\frac1{v+s},\qquad
 \mu\ell=1-\theta,\qquad
 \theta=\frac{s}{v+s},\qquad
 \alpha=\frac{ps_1}{s}=\frac{\theta_+}{\theta}\quad(s>0).
\end{equation}
In particular, the moment in Section~\ref{mc:section} is this law's
actual $\rho$. If its physical support is in $[-R_-,R_+]$, take
\begin{equation}\label{loc:conditional-support}
 a_0=\Delta_-/\ell,\quad a_1=\Delta_+/\ell,\quad
 O_0=R_-/\ell-p,\quad O_1=R_+/\ell-q.
\end{equation}
These give precisely $-O_0\le U_0\le a_0$ and
$-a_1\le U_1\le O_1$. The gap yields $a_0+a_1<1$.
Any zero outer cap can be enlarged to a positive cap when a comparison
formula requires division by it.
There is no probability symmetry in this construction: the original
positive discrepancy determines the orientation of $p$.

For explicit resource coefficients, suppose the already acquired cell
has $p\in I=[l,u]\subset(0,1)$, $\theta\le\Theta<1$,
$A_-\ge A_l>0$, $A_+\ge B_l>0$. Let
\begin{gather*}
 v_* =\max_{p\in I}p(1-p),\quad S_*=A_l+B_l,\\
 b=\frac{(6+2\sigma)(1-\Theta)^{3/2}}{\sqrt{v_*}}
                      \operatorname{dist}(p_*,I)^2,\quad
 g=6(1-\Theta)S_*,\\
 K_-=3(\sigma+3)A_l,\quad K_+=3(\sigma+1)B_l,\quad
 h_{n,\pm}=\frac{K_\pm(1-\Theta)}{nv_*}.
\end{gather*}
Dropping the nonnegative product in \eqref{loc:joint-defect} and
using \eqref{loc:canonical-identities} proves the single inequality
\begin{equation}\label{loc:affine-resource}
 D\ge b+g(a_0+a_1)
       +\lambda_n\{(1-\alpha)h_{n,-}+\alpha h_{n,+}\}.
\end{equation}
It supplies the upstream prices in \eqref{mc:signed-plane}, with
$\mathfrak R=\overline D-b$ whenever
$D\le\overline D$. Replacing $K_\pm$ by their minimum gives the
scalar noise price used by the finite rows. Separate upper bounds for
inward displacement and noise do not imply \eqref{loc:affine-resource}.

The stronger signed prices used in the new comparison are finite
consequences of this inequality. Write $R=\overline D-b\ge0$ and
let $\mu_l>0$ be an established lower bound for $\mu$. The endpoint
floors may first be improved by
\begin{align*}
 A_l&\leftarrow\max\left(A_l,
   \sqrt{\frac{(1-\Theta)l}{1-l}}-\frac{R}{6\mu_lS_*}\right),\\
 B_l&\leftarrow\max\left(B_l,
   \sqrt{\frac{(1-\Theta)(1-u)}u}-\frac{R}{6\mu_lS_*}\right),
\end{align*}
using the old $S_*$ in the denominator. Recompute $g,h_{n,\pm}$ from
these proved floors. Combining \eqref{loc:affine-resource} with
$a_0+a_1\ge s_0/O_0+s_1/O_1$ gives
\[
 \lambda_n\{(1-\alpha)C_0(p)+\alpha C_1(p)\}\le R,
 \quad C_0(p)=h_{n,-}+\frac{g}{n(1-p)O_0},\quad
 C_1(p)=h_{n,+}+\frac{g}{npO_1}.
\]
For fixed positive caps, $C_0$ increases and $C_1$ decreases in $p$.
Thus the uniform allocation plane is
\begin{equation}\label{loc:uniform-allocation}
 \lambda_n\{C_0(l)+[C_1(u)-C_0(l)]\alpha\}\le R.
\end{equation}
When the separately checked slope $C_1(u)-C_0(l)$ is nonnegative,
$\lambda_n\le R/C_0(l)$. More generally, positivity of both
endpoint coefficients always gives
\begin{equation}\label{loc:two-endpoint-noise}
 \lambda_n\le\frac{R}{\min\{C_0(l),C_1(u)\}}.
\end{equation}
This consequence permits either ordering of the signed noise prices.
It also contains the scalar harmonic bound. Indeed, for
$B_* =\max(O_0,O_1)\max(u,1-l)$ and
$h_* =\min(h_{n,-},h_{n,+})$,
\begin{equation}\label{loc:harmonic-consequence}
 \min\{C_0(l),C_1(u)\}\ge h_*+\frac{g}{nB_*},\qquad
 \lambda_n\le\frac{R}{h_*+g/(nB_*)}.
\end{equation}
This follows from $(1-l)O_0\le B_*$ and $uO_1\le B_*$.
It improves $\theta$ through
\eqref{loc:canonical-identities}, and hence improves the next endpoint,
outer-cap and price estimates. A prescribed finite number of these
steps proves the strengthened rows by induction; each step uses the
previous variance bound. It assumes neither a fixed point nor a
special value of the unknown allocation $\alpha$.

The saved floor-$25$ enclosures are interpreted as
\begin{equation}\label{loc:weighted-transfer}
 \mathcal R_n\le r_{25},\qquad nD\le25D_{25}.
\end{equation}
This is a fixed sequence of inequalities for the same actual-$n$ law.
Appendix~\ref{ut:section} gives the complete low/high-frequency
candidates, their index gates and the finite stage induction.
The first unseeded raw bound licenses removal of the Gaussian weight
before its defect output is used. Every subsequent step first
transfers the incoming physical bounds, acquires the needed support
at the actual integer, and then restores the weighted defect bound.
Neither extremality at $25$ nor monotonicity of an adaptively rounded
program is assumed. Physical deductions may use the weaker $D\le D_{25}$.

For the exported plane set $r=25/n$, $K=\min(K_-,K_+)$,
$h_{25}=K(1-\Theta)/(25v_*)$ and $R=D_{25}-b$.
If $R>0$, write $d=R/g$ and $L=R/h_{25}$. Then
the resulting scaled plane is
\begin{equation}\label{loc:scaled-plane}
 \frac{a_0+a_1}{rd}+\frac{\lambda_n}{L}
 \le1-(n/25-1)\frac bR\le1.
\end{equation}
The joint noise intercept $L$ stays fixed. Appendix~\ref{ut:section}
derives this plane, its zero-residual boundary and the weighted
specialization of \eqref{loc:harmonic-consequence}.

Finally, new physical support can improve these same parameters.
From an established $\theta\le t<1$, choose
$0<\ell_l\le\sqrt{(1-t)/v_*}$. Then
\[
 O_0\le R_-/\ell_l-l,\qquad
 O_1\le R_+/\ell_l-(1-u).
\]
Intersect with previous caps and apply the same endpoint elimination;
\eqref{loc:canonical-identities} then updates the variance and mean
bounds. Each step consumes the preceding bound. A scalar consequence
does not replace the joint intercept $L$ in \eqref{loc:scaled-plane}.

Consequently a finite acquisition must cover the entire initial moment
and probability domain by closed cells; prove the support, Fourier and
coefficient inequalities on each cell; and either exclude it by a strict
contradiction or export the simultaneous quantities above. Its surviving
rows are necessary conditions on the original global extremizer.
For the surviving middle-band profiles, the acquired contact radii
satisfy $H_-,H_+<2$, the original $\Lambda<1/2$, and
$p\in[.32,.43]$. Since $\lambda_n\le\Lambda$ and $n\ge26$,
$\sqrt{pq+\lambda_n/n}<3/5$. Hence the raw threshold belongs to
\[
 np+H/\ell\in(np-6/5,np+6/5)\subset(0,n).
\]
At this raw threshold the comparison uses
$np+H/\ell=k+t$ with $|t|\le1/2$, including both choices at a tie.
Every nearest integer lies in $0\le k\le n$; all such integers are
covered or excluded by the acquired contact bounds.
This supplies the continuous meaning of the finite
comparison inputs, independently of their file labels.

% Substantive retained derivation: work/berry-esseen-low26-30-20261026/source/
% cutoff25_20261011/new_root/initial25_b1_exact_v1/RETAINED_ANALYTIC_FRAMEWORK.tex,
% contact/support, identity_fourier, resource_enclosure, and joint-gap sections.
% The fixed-floor interpretation and exact beta0 additions are in source/analytic/
% UNIFORM25_CANTELLI_PRODUCTGAP_REVIEW.txt and
% UNIFORM25_BETA0_SOURCE_AND_COMPOSITION_REVIEW.txt.
% Scaled inward plane: source/original/new_audit/
% SCALED_INWARD_PHYSICAL_DERIVATION_V1.txt.
